\subsection*{用简单函数保留局部信息}
假设要在零附近估算 $\mathrm e^x$，直接计算指数不如多项式方便。先用常数一匹配函数值，再加 $x$ 匹配斜率，再问二次项该是多少。若写 $1+x+cx^2$，它在零处的二阶导数为 $2c$，而指数的二阶导数为一，所以 $c=1/2$。

继续下去，每次求导都会带下一个整数因子，$x^k$ 求 $k$ 次导数后得到 $k!$。所以要匹配原函数的第 $k$ 阶导数，系数必须除以 $k!$。泰勒公式由此来自逐阶匹配，而不是先出现一个无穷求和式。

\subsection*{多加一项不等于宣布全局正确}
在动画里比较不同阶多项式与指数：中心附近先贴合，远处仍可能明显偏离。近似是否够用取决于输入范围和容许误差，不能只说“已经用了三阶”。若在区间上能界住下一阶导数，就能用余项估计把误差写成具体上界。

例如 $\mathrm e^{0.1}$ 的二阶近似为 $1.105$，而余项大小不超过 $\mathrm e^{0.1}0.1^3/6$。这里多出来的三次幂解释了小输入带来的精度；当输入很大时，同样阶数未必有用。

有限阶近似还不等于无穷级数表示。即使函数无穷可微，所有导数信息也可能不足以恢复邻域中的函数；下面的反例提醒我们核对解析性或余项趋零条件。Lean 对照给出一个精确多项式案例，不能替代一般余项定理。
